%! Author = sriadityavedantam
%! Date = 4/8/26

\section{Degree of a Divisor}

Let $X$ be a smooth projective curve.
Let
\[
    D = \sum a_i p_i \in \Div(X).
\]

\begin{definition}[Degree of a Divisor]
    The degree of $D$ is
    \[
        \deg(D) \coloneqq \sum a_i.
    \]
\end{definition}

Let $X = \pp^1$.
A divisor $D$ is principal if and only if $\deg(D) = 0$.
Moreover,
\[
    Cl(X) = \frac{\Div(X)}{\{\text{principal divisors}\}} \cong \z,
\]
via the map
\[
    [D] \mapsto \deg(D).
\]

\begin{theorem*}{
    Let $X$ be a smooth projective curve.
    If $D$ is principal, then $\deg(D) = 0$.
}
\end{theorem*}

\begin{proof}
    Suppose $D$ is principal.
    Then
    \[
        D = \div(f) = \sum \nu_{p_i}(f)\, p_i
    \]
    for some $f \in \c(X)$.
    Thus
    \[
        D = \sum Z(f) - \sum P(f),
    \]
    where $Z(f)$ and $P(f)$ denote the zeros and poles of $f$ respectively.

    It suffices to show that
    \[
        \deg(Z(f)) = \deg(P(f)).
    \]

    If $f$ is nonconstant, then it defines a rational map
    \[
        f : X \dashrightarrow \c.
    \]
    This extends to a morphism
    \[
        f : X \to \pp^1.
    \]
    This map is finite, hence a ramified covering.
    Define $\deg(f)$ as the number of preimages of a general point.

    We claim that $\deg(f)$ equals the number of preimages of any point in $\pp^1$, counted with multiplicity.
    Indeed, for any $p \in \pp^1$ and $y \in f^{-1}(p)$, there exists a neighborhood $U \subset X$ such that
    \[
        f|_U : U \to \c
    \]
    is locally of the form $z \mapsto z^k$ for some $k \in \n$.
    Thus each point contributes with multiplicity.

    Therefore, counting multiplicities,
    \[
        \deg(Z(f)) = \deg(P(f)),
    \]
    so $\deg(D) = 0$.
\end{proof}

The degree is invariant under linear equivalence.
Hence there is a well-defined group homomorphism
\[
    \deg : Cl(X) \to \z.
\]
Define
\[
    Cl^0(X) \coloneqq \ker(\deg).
\]
Then we have a short exact sequence
\[
    0 \to Cl^0(X) \to Cl(X) \xrightarrow{\deg} \z \to 0,
\]
so
\[
    Cl(X) \cong \z \oplus Cl^0(X).
\]

\begin{theorem*}{
    If $X \neq \pp^1$, then $Cl^0(X) \neq 0$.
    Moreover, it is uncountable.
}
\end{theorem*}

In fact,
\[
    Cl^0(X) \cong \mathscr{J}(X),
\]
the Jacobian of $X$.
If $X$ has genus $g$, then $\mathscr{J}(X)$ is a $g$-dimensional abelian variety.
As a complex manifold,
\[
    \mathscr{J}(X) \cong \frac{\c^g}{\z^{2g}}.
\]